\subsection{COMPOSITION OF AN ORDERED SEQUENCE OF ELEMENTS}

Recall that a family of elements of a set E is a mapping \(t \mapsto x_{t}\) of a set I (called an indexing set) into E; a family \((x_{t})_{t \in I}\) is called finite if the indexing set is finite.

A finite family \((x_{t})_{t\in I}\) of elements of E whose indexing set I is totally ordered is called an ordered sequence of elements of E.

In particular, every finite sequence \((x_{i})_{i\in\mathbb{H}}\) , where H is a finite subset of the set N of natural numbers, can be considered as an ordered sequence if H is given the order relation induced by the relation \(m \leqslant n\) between natural numbers.

Two ordered sequences \((x_{i})_{i\in\mathbf{I}}\) and \((y_{k})_{k\in\mathbf{K}}\) are called similar if there exists an ordered set isomorphism \(\phi\) of I onto K such that \(y_{\phi(i)} = x_{i}\) for all \(i \in I\) .

Every ordered sequence \((x_{\alpha})_{\alpha \in \mathbf{A}}\) is similar to a suitable finite sequence. For there exists an increasing bijection of \(\mathbf{A}\) onto an interval \([0, n]\) of \(\mathbf{N}\) .

DEFINITION 4. Let \((x_{\alpha})_{\alpha \in \mathbf{A}}\) be an ordered sequence of elements in a magma \(\mathbf{E}\) whose indexing set \(\mathbf{A}\) is non-empty. The composition (under the law \(\top\) ) of the ordered sequence \((x_{\alpha})_{\alpha \in \mathbf{A}}\) , denoted by \(\bigcap_{\alpha \in \mathbf{A}} x_{\alpha}\) , is the element of \(\mathbf{E}\) defined by induction on the number of elements in \(\mathbf{A}\) as follows:

\begin{enumerate}
\def\labelenumi{(\arabic{enumi})}
\item
  if \(\mathbf{A} = \{\beta\}\) then \(\bigcap_{\alpha \in \mathbf{A}} x_{\alpha} = x_{\beta}\) ;
\item
  if A has \(p > 1\) elements, \(\beta\) is the least element of A and \(A' = A - \{\beta\}\) , then \(\bigcap_{\alpha \in A} x_{\alpha} = x_{\beta} \top \left( \bigcap_{\alpha \in A} x_{\alpha} \right)\) .
\end{enumerate}

It follows immediately (by induction on the number of elements in the indexing sets) that the compositions of two similar ordered sequences are equal; in particular, the composition of any ordered sequence is equal to the composition of a finite sequence. If \(\mathbf{A} = \{\lambda, \mu, \nu\}\) has three elements ( \(\lambda < \mu < \nu\) ) the composition \(\bigcap_{\alpha \in \mathbf{A}} x_{\alpha}\) is \(x_{\lambda} \top (x_{\mu} \top x_{\nu})\) .

Remark. Note that there is a certain arbitrariness about the definition of the composition of an ordered sequence; the induction we introduced proceeds ``from right to left''. If we proceeded ``from left to right'', the composition of the above ordered sequence \((x_{\lambda}, x_{\mu}, x_{\nu})\) would be \((x_{\lambda} \top x_{\mu}) \top x_{\nu}\) .

\[
\left(x _ {\alpha}\right) _ {\alpha \in A}
\]

\[
\begin{array}{c} \underline {{\mathsf {I}}} \\ \alpha \in \mathbf {A} \end{array}
\]

denoted by \(\sum_{\alpha \in \mathbf{A}} x_{\alpha}\) and called the sum of the ordered sequence \((x_{\alpha})_{\alpha \in \mathbf{A}}\) (the \(x_{\alpha}\) being called the terms of the sum); for a law written multiplicatively it is usually denoted by \(\prod_{\alpha \in A} x_{\alpha}\) and called the product of the ordered sequence \((x_{\alpha})\) (the \(x_{\alpha}\) being called the factors of the product).†

When there is no possible confusion over the indexing set (nor over its ordering) it is often dispensed with in the notation for the composition of an ordered sequence and we then write, for example for a law written additively, \(\sum_{\alpha} x_{\alpha}\) instead of \(\sum_{\alpha \in A} x_{\alpha}\) ; similarly for the other notations.

For a law denoted by \(\top\) the composition of a sequence \((x_{i})\) with indexing set a non-empty interval \((p,q)\) of \(\mathbf{N}\) is denoted by \(\bigcap_{p\leqslant i\leqslant q}x_{i}\) or \(\bigcap_{i = p}^{q}x_{i}\) ; similarly for laws denoted by other symbols.

Let E and F be two magmas whose laws of composition are denoted by T and f a homomorphism of E into F. For every ordered sequence \((x_{\alpha})_{\alpha \in \mathbf{A}}\) of elements of E

\[
f \left(\underset {\alpha \in A} {T}\right) = \underset {\alpha \in A} {T} f (x _ {\alpha}).\tag{2}
\]

